\usepackage{fontspec}
\usepackage{newunicodechar}
\usepackage{xcolor}

% XeLaTeX supports native OpenType shaping; HarfBuzz and FallbackFamilies
% are not supported fontspec options for this engine.
\defaultfontfeatures{Ligatures=TeX,Scale=MatchLowercase}
\IfFontExistsTF{Noto Serif}{\setmainfont{Noto Serif}}{\setmainfont{Times New Roman}}
\IfFontExistsTF{Noto Sans}{\setsansfont{Noto Sans}}{\setsansfont{Arial}}
\IfFontExistsTF{Noto Sans Mono}{\setmonofont{Noto Sans Mono}}{\setmonofont{Menlo}}
\usepackage{unicode-math}
\setmathfont{latinmodern-math.otf}

% Emoji and symbols “helper” faces
\IfFontExistsTF{Symbola}{
  \newfontfamily\emoji{Symbola}
}{
  \IfFontExistsTF{Noto Emoji}{
    \newfontfamily\emoji{Noto Emoji}
  }{
    \IfFontExistsTF{DejaVu Sans}{
      \newfontfamily\emoji{DejaVu Sans}
    }{
      \newfontfamily\emoji{Times New Roman}
    }
  }
}
\IfFontExistsTF{Apple Symbols}{
  \newfontfamily\symbolsfont{Apple Symbols}
}{
  \IfFontExistsTF{Symbola}{
    \newfontfamily\symbolsfont{Symbola}
  }{
    \IfFontExistsTF{Noto Sans Symbols}{
      \newfontfamily\symbolsfont{Noto Sans Symbols}
    }{
      \newfontfamily\symbolsfont{Times New Roman}
    }
  }
}

% Emoji wrapper
\newcommand{\textemoji}[1]{%
  {\emoji\strut #1}%
}

% Unicode → robust glyph or math mapping
% Arrows and math-ish
\newunicodechar{→}{{\symbolsfont →}}        % U+2192
\newunicodechar{➝}{{\symbolsfont →}}        % map “long arrow” to same
\newunicodechar{≤}{\ensuremath{\le}}         % U+2264
\newunicodechar{≈}{\ensuremath{\approx}}     % U+2248

% Emojis (add as needed)
\newunicodechar{✅}{\textemoji{✅}}
\newunicodechar{❌}{\textemoji{❌}}
\newunicodechar{🎯}{\textemoji{🎯}}
\newunicodechar{🎬}{\textemoji{🎬}}
\newunicodechar{📹}{\textemoji{📹}}
\newunicodechar{🧠}{\textemoji{🧠}}
\newunicodechar{🔧}{\textemoji{🔧}}
\newunicodechar{📣}{\textemoji{📣}}
\newunicodechar{📈}{\textemoji{📈}}
\newunicodechar{🔍}{\textemoji{🔍}}
\newunicodechar{🔜}{\textemoji{🔜}}
\newunicodechar{🎥}{\textemoji{🎥}}
\newunicodechar{💄}{\textemoji{💄}}
\newunicodechar{🛠}{\textemoji{🛠}}
\newunicodechar{🤝}{\textemoji{🤝}}
\newunicodechar{💭}{\textemoji{💭}}
\newunicodechar{🎲}{\textemoji{🎲}}
\newunicodechar{🌱}{\textemoji{🌱}}
% Non-breaking hyphen (U+2011) → regular hyphen with no break after it
\newunicodechar{^^^^2011}{\nobreakdash-}
% (optional) Hyphen (U+2010 and U+2011) as well, some sources use it
\newunicodechar{^^^^2010}{-}
\newunicodechar{^^^^2011}{\mbox{-}}

% Arrow, works in body and in PDF bookmarks
\providecommand{\texorpdfstring}[2]{#1}
\newcommand{\symbolarrow}{\texorpdfstring{\ensuremath{\rightarrow}}{->}}

% U+2E3B THREE-EM DASH (⸻)
\providecommand{\texorpdfstring}[2]{#1}
\newcommand{\threeemdash}{\texorpdfstring{\rule[0.55ex]{3em}{0.12ex}}{---}}
\newunicodechar{⸻}{\threeemdash{}}

% Fallback, if any raw Unicode arrow survives later filters
\newunicodechar{→}{\symbolarrow{}}
\newunicodechar{➝}{\symbolarrow{}}


% Handle U+202F (narrow no-break space) as non-breaking
\catcode"202F=\active
\begingroup
  \lccode`\~="202F
  \lowercase{\endgroup\def~}{\nobreakspace}

% Strip VS16 if it sneaks through
\catcode`\️=\active
\def️{}

% Emoji command shims used by the Lua filter
\newcommand{\emojicheckmark}{\textemoji{✅}}
\newcommand{\emojicrossmark}{\textemoji{❌}}
\newcommand{\emojitarget}{\textemoji{🎯}}
\newcommand{\emojifilm}{\textemoji{🎬}}
\newcommand{\emojivideo}{\textemoji{📹}}   % also used for 🎥
\newcommand{\emojibrain}{\textemoji{🧠}}
\newcommand{\emojitool}{\textemoji{🔧}}    % also used for 🛠
\newcommand{\emojimegaphone}{\textemoji{📣}}
\newcommand{\emojichart}{\textemoji{📈}}
\newcommand{\emojisearch}{\textemoji{🔍}}
\newcommand{\emojisoon}{\textemoji{🔜}}
\newcommand{\emojimakeup}{\textemoji{💄}}  % fixed typo
\newcommand{\emojihandshake}{\textemoji{🤝}}
\newcommand{\emojithought}{\textemoji{💭}}
\newcommand{\emojidice}{\textemoji{🎲}}
\newcommand{\emojiSeedling}{\textemoji{🌱}}
\newcommand{\emojiGlobe}{\textemoji{🌐}}   % globe with meridians

